\section{Additional Variants}
\label{app:variants}

\subsection{Global variables}
\label{app:globals}

\citet{wu2023probabilistic} introduce global variables (their Appendix~B.3) as an in-graph
substitute for the feed-forward block. Of their three designs we use \emph{single-split}, with
one global variable per position shared by all channels (the others are \emph{all-dep} and
\emph{dep-split}), which they identify as the closest to a transformer's computation.

The construction is position-local, so it carries over unchanged. Slot $t$ gains a variable
$\Gv{t}$ over $\nglob$ values and one binary factor
$\phi_b(\Gv{t}{=}k, \Zv{t}{=}a) = \exp(\Bglob_{k,a})$ with
$\Bglob \in \R^{\nglob \times \nlab}$, so \Cref{eq:slot-energy} gains the term
$-\sum_{k,a} \QG{t}(k)\,\QZ{t}(a)\,\Bglob_{k,a}$. There is no channel axis, so the message
enters once per iteration. The slot is still a tree ($\Gv{t}$ is one more leaf off $\Zv{t}$), so
\Cref{prop:exact-readout} still applies; $\Gv{t}$ reads no other position, so causality is
untouched; and the cost is $O(\nglob\nlab)$ per position. The mean-field updates gain one pair
of messages,
%
\begin{equation}
\label{eq:global-update}
\begin{aligned}
  \QG{t}(k) &\propto \exp\Big(\tfrac{1}{\lamG}\textstyle\sum_a \QZ{t}(a)\,\Bglob_{k,a}\Big),\\
  \mathcal{G}^{\mathrm{glob}}_{t}(a) &= \textstyle\sum_k \QG{t}(k)\,\Bglob_{k,a},
\end{aligned}
\end{equation}
%
the second of which joins the arc messages inside the bracket of \Cref{eq:update-Z}. Together
they form the operator $\sigma(q\Bglob^{\!\top}\!/\lamG)\Bglob$ on the label belief $q$, which
Wu and Tu identify as the analogue of a feed-forward block: a map into $\nglob$ features, a
softmax in place of a \textsc{relu}, and a map back through the transpose of the same matrix.
Wu and Tu write this operator without a message weight; we carry $\lamG$, and $\lamG = 1$
recovers their form.

\paragraph{Under the exact readout the global variable acts only as a prior.}
Folding $\Gv{t}$ into the sum--product of \Cref{prop:exact-readout} multiplies $\readout{t}(a)$
by $\sum_k e^{\Bglob_{k,a}}$:
%
\begin{equation}
\label{eq:global-readout}
  \log \readout{t}(a) = \textstyle\sum_{c} \LSE_{j \in \Dep{t}} \Bfac{c}_{j,a}
    \;+\; \LSE_{k} \Bglob_{k,a}.
\end{equation}
%
The added term is one $\nlab$-vector, independent of $t$ and of the prefix: a fixed per-label
tilt of the mixture weights. It does not cancel, since the readout takes a log-sum-exp over
labels, but it cannot carry context. The variable still reshapes every $\qbar{i}$ through the
observed-mode beliefs, so under the exact readout, which has no predictive inner loop, a run would mix
a constant tilt with an indirect effect through the prefix. Our implementation refuses this
combination, and validation checks assert that the direct term is constant and the end-to-end
effect is not (\Cref{app:validation}). \Cref{sec:exp-ablations} measures the mean-field
configuration, in which the variable also acts at the predicted slot.

\subsection{Serial update schedule}
\label{app:serial}

Both schedules are causal. The \emph{serial} schedule runs \Cref{eq:update-H,eq:update-Z} at
position $i$ for a fixed number of inner rounds, freezes $\qbar{i}$ and moves right. Each slot
is then a well-posed mean-field problem on its own energy, with the per-step guarantees of the
entropic Frank--Wolfe scheme, and causality follows by induction on $i$. The
\emph{layer-parallel} schedule computes all $\QZ{i}^{(t)}$ from $\{\QZ{j}^{(t-1)}\}_{j<i}$ at
once under a strict lower-triangular mask, for $\niter$ iterations. Its computation graph is
that of a depth-$\niter$ causal transformer with shared weights. All reported numbers use this
schedule, because only under it do the causal PT and the looped transformer share a computation
graph and an effective depth, so that their difference isolates structure.

The two schedules do not compute the same thing. We compared their observed-mode label beliefs at $\nlab=32$, $\nchan=4$, $\krank=8$, $\clip=3$, $\niter=5$ and length $24$. The mean
total-variation distance is $3.4\times10^{-6}$ at initialisation scale $0.02$, $0.111$ at
$0.5$ and $0.682$ at $2.0$, where the maximum reaches $1.000$ (disjoint label mass at some
positions). On an overfitted model it falls below $10^{-4}$. The schedules agree when beliefs
are near uniform or near one-hot, and can disagree completely in between.

\subsection{Root node}
\label{app:root}

The root is not a variable in this model. The head domain is
$\Dep{t} = \{\textsc{root},1,\dots,t-1\}$, and $\rfac{c} = \Bfac{c}_{\textsc{root},\cdot}$ is a
learned column of $\nlab$ numbers per channel. Like a beginning-of-sequence token, it gives
position $1$, whose head domain would otherwise be empty, something to attend to. The root of
\citet{wu2023probabilistic} (their \S2.3.5) is an optional latent variable with its own label set
and ternary potential. Ours is a constant parameter: it has no posterior, carries no information
between positions and so cannot break causality, and costs $\nchan\nlab$ numbers instead of a
$\nlab \times \nlab_{\textsc{root}} \times \nchan$ tensor. It is their root variable with its
distribution held fixed and contracted into a binary factor as in \Cref{eq:contract}, the
collapse argument Wu and Tu use for global feature variables (their Appendix~B.3.1).

At initialisation $\rfac{c}$ enters the attention in raw $\nlab$-space, while the prefix rows
arrive contracted, $\Bfac{c}_{j,a} = \E_{\qbar{j}}[\Tfac{c}_{a,\cdot}]$, shrunk by roughly
$\nlab^{-1/2}$ for a near-uniform belief and again by the low-rank factorisation
$\Tfac{c}=\Ufac{c}\Vfac{c}{}^{\top}$. With all parameters
drawn from one $\mathcal{N}(0,\sigma^2)$, the \textsc{root} row has $121\times$ the norm of the
prefix rows at $\nlab=384$, $\nchan=16$, $\krank=64$. A near-uniform $\QZ{t}$ hides this at
first, but it could draw most of the attention once beliefs sharpen. In our measurements it
does not: untrained, in the configuration above, the root receives $1.10\times$ uniform attention
mass at the last slot, and on an overfitted toy model it receives $0.0001$ against $0.125$ under
uniform attention. How much mass a trained language model sends to the root still depends on the
configuration (\Cref{app:interpretability}). The initialisation scale of $\rfac{c}$ is a separate
hyperparameter, left equal to that of the other parameters.

\subsection{Factored labels}
\label{app:factored}

Factored labels replace the label variable at each position by $\ncomp$ variables. As
elsewhere, $\nlab$ is the total label width, so each component ranges over
$\nlabc = \nlab/\ncomp$ values. Position $t$ carries $\Zv{t}^{(1)},\dots,\Zv{t}^{(\ncomp)}$ and
$\ncomp$ word--label factors that share the word variable,
$\phi\big(\Wv{t},\Zv{t}^{(k)}\big) = \exp\big(\Sfac^{(k)}_{\Wv{t},\Zv{t}^{(k)}}\big)$. Each
channel $c$ has a child component $k(c)$, whose belief queries the channel and receives its
message, and a head component $k'(c)$, whose prefix belief forms the channel's keys:
$\Bfac{c}_{j,a} = \sum_b \qbar{j}^{(k'(c))}(b)\,\Tfac{c}_{a,b}$. The per-slot energy becomes
%
\begin{equation}
\label{eq:factored-energy}
\begin{aligned}
  E_t = &-\textstyle\sum_w \QW{t}(w)\,\bfac_w \\
        &-\textstyle\sum_k \sum_{w,a} \QW{t}(w)\,\QZ{t}^{(k)}(a)\,\Sfac^{(k)}_{w,a}\\
        &-\textstyle\sum_c \sum_{j\in\Dep{t}} \sum_a \QH{t}{c}(j)\,\QZ{t}^{(k(c))}(a)\,
           \Bfac{c}_{j,a},
\end{aligned}
\end{equation}
%
still multilinear, so the updates are still softmaxes of the gradients. Causality is untouched,
the attention correspondence holds per channel with component-specific queries and keys, and the
compute is linear: $\ncomp$ components of width $\nlabc$ cost what one of width $\nlab$ costs,
since each channel touches one component. The slot
graph is still a tree ($\Wv{t}$ joined to $\ncomp$ stars), so \Cref{eq:exact-readout} becomes a
product of $\ncomp$ mixtures,
%
\begin{equation}
\label{eq:factored-readout}
\begin{aligned}
  \hat p(w) &\propto e^{\bfac_w}\prod_{k=1}^{\ncomp}
     \sum_{a=1}^{\nlabc} e^{\Sfac^{(k)}_{w,a}}\,\readout{t}^{(k)}(a),\\
  \log \readout{t}^{(k)}(a) &= \sum_{c\,:\,k(c)=k} \LSE_{j\in\Dep{t}} \Bfac{c}_{j,a},
\end{aligned}
\end{equation}
%
a strictly richer family than one mixture at the same total width.

Under the mean-field readout, factoring adds no expressivity over a flat label of width
$\nlab$. The
logits $\bfac + \sum_k \Sfac^{(k)}q^{(k)}$ are affine in the concatenated beliefs, so their
affine dimension is at most $\nlab$, the bound of \Cref{prop:readout}(i) for $\ncomp=1$. Any gain
in the head must come through \Cref{eq:factored-readout}. Factoring also changes the inference
dynamics and the parameter count, so \Cref{sec:exp-factored} tests the interaction: the
prediction is an effect of $\ncomp$ under the exact readout that is absent under the mean-field
readout, not an effect of the same size under both.

\paragraph{Channel assignment.}
We use $k(c) = c \bmod \ncomp$ and $k'(c) = \lfloor c/\ncomp\rfloor \bmod \ncomp$, so that at
$\nchan = \ncomp^2$ every ordered pair of components occurs exactly once. The alternative
$k'(c) = k(c)$, in which components never read each other, would separate the effect of having
more mixture components from the effect of letting components communicate. We do not combine global variables with factored labels.

\paragraph{Parameter count.}
The word--label factors keep $|\Vocab|\,\ncomp\,\nlabc = |\Vocab|\,\nlab$ parameters, as for a
flat label of width $\nlab$; $\Sfac$ stays one tensor, read as $\ncomp$ blocks of width
$\nlabc$. The arc factors, however, act on $\nlabc\times\nlabc$. The non-embedding block,
$2(\clip{+}1)\nchan\,\nlabc\,\krank$ for the arc scores plus $\nchan\,\nlabc$ for the root
column, therefore falls as $1/\ncomp^2$ when the rank scales with $\nlabc$ and as $1/\ncomp$
when it is fixed. At $\nchan=4$, $\clip=3$ and $\krank=\nlabc$ the measured non-embedding counts
are $32{,}896$ at $(\nlab,\ncomp) = (32,1)$, $8{,}256$ at $(32,2)$, $32{,}896$ at $(64,2)$,
$8{,}256$ at $(64,4)$ and $32{,}896$ at $(128,4)$. Factoring therefore never gains through extra
parameters; at a fixed non-embedding budget it gains total label width: the three equal counts
are models of one size at $\nlab = 32$, $64$ and $128$.
